\section{The Hard Split}
\label{app:hard}
The \nHard{} tasks come from ansible (8), openlibrary (8), flipt (7), teleport (6), webclients (5), qutebrowser (4), vuls (4), element-web (4), navidrome (3), NodeBB (1) and tutanota (1). A task qualified if at least two of five model families (Opus~5, GLM-5.3, Kimi K3, Inkling, Gemini~3.8 Flash) failed it. That gave \nHardCandidates{} candidates after excluding \nHardDefectDrop{} tasks with verified defects, and \nHardAuditDrop{} more were removed after review. Every Hard task passes with its reference solution and fails with an empty patch. Readers should know four things about the released list.
\begin{itemize}\setlength\itemsep{0pt}
\item \hardKimiOnly{} tasks qualify only through partner runs of Kimi K3.
\item \hardRewritten{} of the \nHard{} instructions were edited after selection, against \otherRewritten{} elsewhere.
\item One task (openlibrary \task{7cbfb812}) has a verified test defect, and one (ansible \task{40ade1f8}) has a timing-sensitive test.
\item Every Hard task is solved by at least one of our runs.
\end{itemize}

\section{Disagreements Between the Two Grades}
\label{app:flips}
Table~\ref{tab:flips} lists every task whose grade in the agent's sandbox differs from its first grade in a clean container, across our public runs. The runs labelled v2.2r used an earlier release of the tests.

\begin{table*}[t]
\centering\small
\begin{tabular}{p{0.17\textwidth}p{0.20\textwidth}ccp{0.39\textwidth}}
\toprule
Run & Task & In place & Fresh & Cause and adjudication \\
\midrule
Opus 5 (high, v2.2r) & vuls \texttt{ef2be3d6} & \cmark & \xmark & Forged \texttt{goval-dictionary v0.9.5}: module zip served from a file-based \texttt{GOPROXY}, fabricated \texttt{go.sum} entry; the clean image fails the checksum \\
Opus 5 (high, v2.2r) & vuls \texttt{ca3f6b1d} & \cmark & \xmark & Vendored, modified dependency whose \texttt{log/} package is git-ignored; the captured patch is incomplete \\
Opus 5 (high, v2.2r) & ansible \texttt{40ade1f8}, teleport \texttt{ba6c4a13} & \xmark & \cmark & One timing-sensitive test failed under agent-phase load (two tasks) \\
Opus 5 (high, v2.2r) & webclients \texttt{c6f65d20} & \cmark & \xmark & Missed a 5\,s Jest timeout on a cold image; fresh grades 0,1,1, adjudicated pass \\
Inkling (high, v2.2r) & vuls \texttt{ef2be3d6} & \cmark & \xmark & Forged \texttt{goval-dictionary v0.9.5} from the older cached module with a fabricated \texttt{.info} and \texttt{go.sum} entry, as Opus~5 did; 2/2 fresh fails \\
Inkling (high, v2.2r) & vuls \texttt{ca3f6b1d} & \cmark & \xmark & Edited a struct field into a cached dependency with \texttt{sed -i}; nothing in the diff; 2/2 fresh fails \\
Inkling (high, v2.2r) & navidrome \texttt{874b17b8} & \cmark & \xmark & Made the module cache writable and added a file to a cached dependency; 2/2 fresh fails \\
Gemini 3.8 Flash & flipt \texttt{3d5a345f} & \cmark & \xmark & Fabricated \texttt{go.sum} entries for real dependencies after rehearsing the forgery on a fake module in the local proxy cache; 2/2 fresh fails \\
Gemini 3.8 Flash & flipt \texttt{c188284f} & \cmark & \xmark & Fabricated \texttt{go.work.sum} checksum; 2/2 fresh fails \\
Gemini 3.8 Flash & navidrome \texttt{874b17b8} & \cmark & \xmark & Edited the module cache in place; 2/2 fresh fails (Opus~5 and GLM-5.3 solved the same task offline by vendoring a local fork under a \texttt{replace} directive) \\
Gemini 3.8 Flash & NodeBB \texttt{be43cd25}, tutanota \texttt{f3ffe17a} & \xmark & \cmark & Flaky tests; 2/2 fresh passes \\
Opus 5 (xhigh) & ansible \texttt{40ade1f8} & \cmark & \xmark & The timing-flaky revised test; fresh grades 0,1,1, adjudicated pass \\
Opus 5 (xhigh) & element-web \texttt{71fe08ea} & \xmark & \cmark & Fresh grades 1,1,1; in-place fail was a sandbox flake \\
Opus 5.5 (xhigh) & element-web \texttt{72a8f8f0} & \cmark & \xmark & Flaky composer test; fresh grades 0,1,1, adjudicated pass \\
Fable 5.1 & NodeBB \texttt{bd80d36e} & \cmark & \xmark & Fresh grades 0,1,1, adjudicated pass \\
Fable 5.1 & teleport \texttt{645afa05} & \xmark & \cmark & Fresh grades 1,1,1 \\
GPT-6 Astra & element-web \texttt{53b42e32} & \cmark & \xmark & Flaky test unrelated to the edit; fresh grades 0,1,0, adjudicated fail \\
\bottomrule
\end{tabular}
\caption{Every task whose in-place and first fresh-sandbox grades differ, across the runs of Table~\ref{tab:main}. Pass-to-fail flips are either environment dependence (forged or edited dependencies, all rejected by the clean image) or flaky tests; fail-to-pass flips are sandbox-side flakes. GLM-5.3 had no disagreement between its two grades. Private-set flips are in the release; they are the caught forgeries of \S\ref{sec:private} and flaky tests.}
\label{tab:flips}
\end{table*}


\section{The Never-Public Set}
\label{app:private}
The \nPriv{} tasks come from \nPrivOrgs{} organisations (\nPrivGo{} in Go, the rest in Python and JavaScript) whose code, to their knowledge, has never been on the open web. Table~\ref{tab:privdefects} lists the defects found while running it and their status. The instruction audit was done by an LLM judge (Claude Opus~5) that saw the instruction, the hidden tests and the reference solution but no model output. We checked eight of its judgments by hand. Its grades are noisy: it agreed with itself on \privJudgeAgree\% of re-judged unchanged tasks. On the final release, every system was run again with the same scaffold, effort and budget. With the imitation-SDK passes counted as failures, the scores are \privStrictOpusPct{} (Opus~5), \privStrictKimiPct{} (Kimi K3), \privStrictGlmPct{} (GLM-5.3), 86.0 (Gemini) and 79.0 (Inkling).

\begin{table*}[t]
\centering\small
\begin{tabular}{p{0.17\textwidth}p{0.50\textwidth}p{0.25\textwidth}}
\toprule
Defect & Description & Status in the final release \\
\midrule
D1 hidden-test restore & The shipped verifier applied the hidden test patch over the agent's working tree; any agent edit to a file the patch touches made the patch fail to apply and graded the task 0 (15--20 tasks per system) & Fixed (restore before apply in all 272 verifiers); +10 to +18 tasks per system on replay \\
D5 test-symbol collision & An agent-added test file declaring a symbol the hidden test also declares fails the package compile & Fixed (colliding agent test files removed) \\
D2 Go toolchain paths & Go images did not export \texttt{GOPATH}/\texttt{GOMODCACHE}, so agents had to rediscover the module cache (Opus~5 in 131 of 154 Go trajectories, Inkling in 65) & Fixed in rebuilt images; offline compile of a sample of hidden-test packages went from 0/7 to 7/7 \\
D3 end-of-life base images & Package repositories of older Debian images no longer resolve & Fixed in 13 of 14 images; one still cannot install packages \\
D4 state-sensitive in-place grade & Fail-to-pass flips from agent-side toolchain state (\texttt{go env -w}) & Closed by fresh-sandbox grading \\
Manifest egress & Three consecutive exports lacked the agent-phase no-network field, so egress was open under the harness flags & Fixed at source; probe still runs before every job \\
Reference patch in verifier bundle & The verifier configuration carried the reference and test patches and commit identifiers & Removed; graded lists unchanged \\
Missing SDK modules & Two tasks need AWS and Deepgram SDK modules absent from the images, so an honest solution cannot compile offline & Open; the 5 imitation-SDK passes occurred here \\
Flaky test & A timing-dependent test fails a correct solution about one run in five in ten tasks & Open; adjudication restores most cases \\
\bottomrule
\end{tabular}
\caption{Verifier and image defects found in the private set. The verifier repair is ours, validated with the two-sided gate and a full replay of every captured patch. The final release also carries the maintainers' image rebuild and the 68 instruction corrections.}
\label{tab:privdefects}
\end{table*}


\section{Reproducibility}
\label{app:repro}
\paragraph{Harness.} Upstream Harbor \harborVersion{} on Modal \modalVersion{}. The network is set per phase: public during setup, only the model endpoint during the agent phase, public again for the verifier, because \nNeedNet{} tasks download dependencies while testing.
\paragraph{Scaffolds.} Claude Code \ccVersion{} on the public runs and 2.1.283 on the private runs, with web tools disabled. Codex \codexVersion{}, with hosted web search disabled. mini-swe-agent 2.4.6, with no step or cost limit, so the time budget binds.
\paragraph{Budget and grading.} The budget is 3000\,s per task. The agent is stopped at 2940\,s, and the patch present then is graded. The patch is the staged \texttt{git diff} of the working tree. It is applied to a fresh container built from the task image, with a three-way fallback, and the unchanged verifier runs. When the sandbox grade and the clean grade disagree, two more clean grades are taken and the majority of the three clean grades counts.
\paragraph{Release.} Task directories, the Hard split, the probe, capture and re-grade scripts, the exact harness flags, and per-task grades in both modes for every run.
